\subsection*{(1) Continuous and strictly increasing transformations}

Let $f: \mathbb{R} \rightarrow \mathbb{R}$ be a continuous and strictly increasing function. We want to show that $VaR_{\alpha}(f(L)) = f(VaR_{\alpha}(L))$.

First, recall the definition of Value at Risk:
\begin{equation*}
VaR_{\alpha}(L) = \inf\{x \in \mathbb{R} : \mathbb{P}(L \le x) \ge \alpha\}
\end{equation*}

Now, let us find the cumulative distribution function for $f(L)$. Since $f$ is strictly increasing, it is invertible, and the event $\{f(L) \le x\}$ is equivalent to the event $\{L \le f^{-1}(x)\}$.
\begin{equation*}
\mathbb{P}(f(L) \le x) = \mathbb{P}(L \le f^{-1}(x)) = F_L(f^{-1}(x))
\end{equation*}

Applying the definition of VaR to $f(L)$:
\begin{equation*}
VaR_{\alpha}(f(L)) = \inf\{x \in \mathbb{R} : \mathbb{P}(f(L) \le x) \ge \alpha\}
\end{equation*}
\begin{equation*}
VaR_{\alpha}(f(L)) = \inf\{x \in \mathbb{R} : F_L(f^{-1}(x)) \ge \alpha\}
\end{equation*}

Let us perform a change of variable by setting $y = f^{-1}(x)$, which implies $x = f(y)$. Because $f$ is strictly increasing, as $x$ increases, $y$ also increases. Furthermore, taking the infimum over $x$ is equivalent to taking the infimum over $y$ and then applying $f$, due to the fact that $f$ is continuous and preserves the ordering of the set.
\begin{align*}
VaR_{\alpha}(f(L)) &= \inf\{f(y) \in \mathbb{R} : F_L(y) \ge \alpha\} \\
&= f\left( \inf\{y \in \mathbb{R} : F_L(y) \ge \alpha\} \right) \\
&= f(VaR_{\alpha}(L))
\end{align*}


\subsection*{(2) Continuous and strictly decreasing transformations}

Let $f: \mathbb{R} \rightarrow \mathbb{R}$ be a continuous and strictly decreasing function. We want to show that $VaR_{\alpha}(f(L)) = f(VaR_{1-\alpha}^{+}(L))$.

Since $f$ is strictly decreasing, the inequality direction reverses when applying the inverse function. The event $\{f(L) \le x\}$ is equivalent to $\{L \ge f^{-1}(x)\}$.
\begin{equation*}
\mathbb{P}(f(L) \le x) = \mathbb{P}(L \ge f^{-1}(x))
\end{equation*}

We write the VaR for $f(L)$:
\begin{equation*}
VaR_{\alpha}(f(L)) = \inf\{x \in \mathbb{R} : \mathbb{P}(L \ge f^{-1}(x)) \ge \alpha\}
\end{equation*}

Let $y = f^{-1}(x)$. Since $f$ is strictly decreasing, as $x$ decreases toward its infimum, $y$ increases toward its supremum. Therefore, finding the infimum of $x$ is equivalent to finding the supremum of $y$ and applying $f$:
\begin{equation*}
VaR_{\alpha}(f(L)) = f\left( \sup\{y \in \mathbb{R} : \mathbb{P}(L \ge y) \ge \alpha\} \right)
\end{equation*}

We must now show that $\sup\{y \in \mathbb{R} : \mathbb{P}(L \ge y) \ge \alpha\} = VaR_{1-\alpha}^{+}(L)$.
Recall the definition of $VaR_{1-\alpha}^{+}(L)$:
\begin{equation*}
VaR_{1-\alpha}^{+}(L) = \inf\{z \in \mathbb{R} : \mathbb{P}(L \le z) > 1 - \alpha\}
\end{equation*}

Let $A = \{y \in \mathbb{R} : \mathbb{P}(L \ge y) \ge \alpha\}$ and $B = \{z \in \mathbb{R} : \mathbb{P}(L \le z) > 1 - \alpha\}$.
If $y \in A$, then $\mathbb{P}(L \ge y) \ge \alpha$. Since $\mathbb{P}(L \ge y) = 1 - \mathbb{P}(L < y)$, we have:
\begin{align*}
1 - \mathbb{P}(L < y) &\ge \alpha \\
\mathbb{P}(L < y) &\le 1 - \alpha
\end{align*}
For any $z < y$, it is strictly true that $\mathbb{P}(L \le z) \le \mathbb{P}(L < y) \le 1 - \alpha$. This implies $z \notin B$. Therefore, any element in $B$ must be greater than or equal to $y$. This means $y \le \inf B$, which implies $\sup A \le \inf B$.

Conversely, take any $y < \inf B$. By definition of the infimum of $B$, $y \notin B$, which means $\mathbb{P}(L \le y) \le 1 - \alpha$. 
Consequently, $\mathbb{P}(L > y) = 1 - \mathbb{P}(L \le y) \ge \alpha$. 
Since $\mathbb{P}(L \ge y) \ge \mathbb{P}(L > y) \ge \alpha$, it follows that $y \in A$. 
Since every value less than $\inf B$ is in $A$, we must have $\sup A \ge \inf B$.

Combining $\sup A \le \inf B$ and $\sup A \ge \inf B$, we conclude $\sup A = \inf B$. 
Thus:
\begin{equation*}
\sup\{y \in \mathbb{R} : \mathbb{P}(L \ge y) \ge \alpha\} = \inf\{z \in \mathbb{R} : \mathbb{P}(L \le z) > 1 - \alpha\} = VaR_{1-\alpha}^{+}(L)
\end{equation*}

Substitute this back into the equation for $VaR_{\alpha}(f(L))$:
\begin{equation*}
VaR_{\alpha}(f(L)) = f(VaR_{1-\alpha}^{+}(L))
\end{equation*}

\vspace{0.5cm}
\noindent \textbf{Deduction for continuous $F_L$:}

We are given the property that if $F_L$ is continuous, $VaR_{\beta}(L) = VaR_{\beta}^{+}(L)$ for any confidence level $\beta \in (0,1)$.
Applying this property to $\beta = 1 - \alpha$, we get:
\begin{equation*}
VaR_{1-\alpha}^{+}(L) = VaR_{1-\alpha}(L)
\end{equation*}

Substituting this into the result we just proved for strictly decreasing functions:
\begin{equation*}
VaR_{\alpha}(f(L)) = f(VaR_{1-\alpha}(L))
\end{equation*}